\clearpage\SourceLeaf{526}\pdfbookmark[1]{Für Ihr gütiges Schreiben}{gauss-dirichlet-fur-ihr-gutiges-schreiben}
Für Ihren freundlichen Brief und die freundliche Übersendung Ihrer beiden Abhandlungen danke ich Ihnen, mein sehr geschätzter Freund, aufs Verbindlichste. Mit Freude sehe ich das wachsende Interesse, das man jetzt den Untersuchungen der höheren Arithmetik entgegenzubringen beginnt. Die gelungene Art, wie Sie das Zweite über die biquadratische Residualität der Zahl 2 aus dem Ersten ableiten, hat mir sehr gefallen.
\par Vermutlich hat der 6. Band unserer Commentationen jetzt Breslau erreicht, sodass Ihnen meine Commentatio prima über die biquadratischen Reste inzwischen wohl bekannt ist. Bei einer passenden Gelegenheit würde ich Ihnen auch gern einen Sonderabdruck übersenden. Für das darin vorkommende Theorem hätte ich zwischen mehreren Beweisarten wählen können; Ihnen wird aber nicht entgehen, warum ich hier den dort ausgeführten Beweis vorgezogen habe: vor allem deshalb, weil die Einordnung von 2 bei den Moduln, für die es quadratischer Nichtrest ist (unter $B$ oder $D$), als wesentlicher Bestandteil des Theorems gelten muss, auf den die meisten anderen Beweisarten nicht anwendbar scheinen.
\par Die ganze Untersuchung, deren Stoff ich schon seit 23 Jahren vollständig besitze, die Beweise der Haupttheoreme (zu denen das Theorem in der ersten Abhandlung noch \textit{nicht} zu rechnen ist) aber seit etwa 14 Jahren \textemdash{} (wobei ich wünsche und hoffe, an Letzteren, den Beweisen, noch einiges vereinfachen zu können) \textemdash{}, habe ich auf ungefähr 3 Abhandlungen angelegt. Mit der Abfassung der zweiten habe ich bereits begonnen und hoffe, sie bald zu vollenden, sofern die mir neuerdings wieder übertragenen Vermessungsarbeiten dabei nicht noch eine Verzögerung verursachen.
\par Das Schlusstheorem $b\equiv\frac12rr\, (\mathrm{mod}.p)$ hatte ich schon vor drei Jahren ebenfalls in den hiesigen gel. Anzeigen bekannt gemacht und auf den merkwürdigen Knoten hingewiesen, der dabei noch zu lösen war; bisher habe ich aber nicht gehört, dass jemand einen Versuch dazu unternommen hätte. Vor einigen Tagen ist mir nun die Lösung für die \textit{eine Hälfte} tatsächlich gelungen. Dieser Fund hat mir umso mehr Freude gemacht, als er sich überhaupt nicht auf Induktion stützt \textemdash{} denn ich gestehe, dass ich gerade \textit{diesen} Zusammenhang nicht erwartet hätte \textemdash{}, sondern a priori auf die Verbindung anderer \textit{sehr verschlungener} und interessanter, schon 28 Jahre alter, aber noch völlig unveröffentlichter Untersuchungen,
